\section{introduction}
One of the most important goals of QCD is to understand the hadron structure from its fundamental degrees of freedom - quarks and gluons. This necessarily goes to the nonperturbative regime of QCD, and it is difficult in general to obtain first principle results directly from the QCD Lagrangian. A powerful tool of obtaining such results is Lattice QCD, which is an approach defined on Euclidean spacetime, and has been used to calculate hadron masses, charges etc. to a remarkable accuracy. However, it cannot be used to directly access intrinsically Minkowskian quantities such as the parton distribution functions (PDFs). PDFs are defined as the forward hadronic matrix elements of light-cone correlations, and describe the momentum distribution of quarks and gluons inside the hadron. They play a crucial role in understanding the experimental data at high energy hadron colliders such as the LHC. Owing to their real time dependence, Lattice QCD can only be used to indirectly extract the information on the PDFs by calculating their moments, which is limited by technical complications such as operator mixing. Another commonly used approach to determine PDFs is to assume a suitable parametrized form and fit to a large variety of experimental data. Most PDF groups obtained their PDF sets in this way~\cite{Alekhin:2012ig,Gao:2013xoa,Radescu:2010zz,CooperSarkar:2011aa,Martin:2009iq,Ball:2012cx}. A main drawback of this approach is that it suffers from parametrization uncertainty, and different groups usually produce different results for the same PDF.


Recent developments~\cite{Ji:2013fga,Ji:2013dva,Xiong:2013bka,Lin:2014zya,Hatta:2013gta,Ma:2014jla,Ji:2014hxa,Ji:2014gla,Ji:2014lra,
Alexandrou:2015rja,Ji:2015jwa,Ji:2015qla,Xiong:2015nua,Li:2016amo,Chen:2016utp} showed that the light-cone observables such as the PDFs can be directly extracted from the large momentum limit of the hadronic matrix element of a spacelike correlator, which is known as the quasi observable, using a large momentum effective theory~\cite{Ji:2014lra} (for other approaches to extract light-cone quantities see e.g.~\cite{Davoudi:2012ya,Detmold:2005gg,Liu:1993cv,Liu:1998um,Liu:1999ak,Liu:2016djw,Braun:2007wv}). The quasi PDF does not have a real time dependence, and thus can be simulated on the lattice. The infrared (IR) behaviors between the flavor non-singlet quasi-PDF and its corresponding light-cone PDF are shown to be the same at one loop by direct computations, and argued to be the same at all loops in Ref.~\cite{Xiong:2013bka}, based on which a factorization formula was also presented.

The factorization in Ref.~\cite{Xiong:2013bka} was given for the bare quasi quark distribution, where all fields and couplings entering the quasi distribution are bare ones. However, as the light-cone distribution, the quasi distribution also contains ultraviolet (UV) divergences and therefore needs renormalization. Ref.~\cite{Ji:2015jwa} explores the renormalization property of the quasi distribution, and shows that it is multiplicatively renormalizable at two-loop order. Also an equivalence was established between the virtual correction of the quasi quark distribution and the correction to the heavy-light quark vector current in heavy quark effective theory, so that the UV divergences in the former can be renormalized as the renormalization of the latter. Dimensional regularization was used in Ref.~\cite{Ji:2015jwa}, and the linear divergence present in a cutoff or lattice regularization was ignored. In realistic lattice calculations, one needs to know how to deal with such power divergences. This is one goal of the present paper. We will show that the power divergence in the quasi quark distribution can be removed to all-loop orders by a mass counterterm, which is the same as the renormalization of an open Wilson line. After such a mass renormalization, the quasi quark distribution is improved such that it contains at most logarithmic divergences.

The rest of this paper is organized as follows. In Sec. 2, we discuss the renormalization of power divergences arising from the Wilson line self energy in the quasi quark distribution. In Sec. 3, we present a lattice perturbation theory matching for the quasi quark distribution. We then conclude in Sec. 4.


